\subsubsection{Data warehouse}
Para el almacenamiento de datos, se optó por usar MongoDB \cite{mongodb}, dada
su amplio uso en el campo de \textit{big data}, su facilidad para almacenar
grandes cantidades de datos de manera distribuída, y el hecho de que es una base
de datos no relacional.

Se creó una única base de datos (\textit{jose}) con dos colecciones: una para
los avistamientos (\textit{sightings}) y otra para los eventos astronómicos
(\textit{astronomical-events}).

Se decidió usar la versión de contenedor de Docker dado que permite una mayor
flexibilidad en el despliegue, al mismo tiempo que está empaquetada con el resto
de la aplicación (para más fácil distribución) y no requiere de instalación en
el sistema.
También se habilitó la persistencia de datos mediante un volumen de Docker, para
así mantener los datos ya extraídos aunque se reinicie la aplicación.